\chapter{对抗安全}
\label{chap:trust-security}

一封钓鱼邮件被误判，可能是语言随时间变化，也可能是发送者反复改写后专门选中了能绕过过滤的表达。后一种情形中，输入由攻击者按目标选择。邮件分类器若还连接归档、检索或发送工具，操纵一次预测又可能影响后续行动，造成信息外泄或越权操作。我们需要判断的因而是：在攻击者能够看见、修改和尝试的范围内，整条处理链能否守住安全要求。

第~\ref{chap:trust-foundations}章的三层架构为这一判断提供前提。系统属性或规范要求层确定受保护的任务目标、隐私要求和行为边界；证据能力层提供测试、机制分析与审计。本章在失效、处置与生命周期层把攻击者的知识、权限、资源和目标写进条件，再沿损害传播的位置安排防御、监测与恢复。

\term{对抗安全}（adversarial security）研究机器学习系统如何抵御有目的的操纵。我们先用攻击的定义、实例和损害明确保护对象，再检查输入、训练过程、模型与运行组件的供应链、模型内部、查询接口及执行链中的薄弱环节；防御据此限制攻击机会、模型受影响的程度和已经发生的损失。分布偏移（第~\ref{sec:distribution-shift}节）描述训练与评价分布的差异，算法稳定性（第~\ref{sec:algorithmic-stability}节）比较训练样本变化后的输出或损失，隐私泄露（第~\ref{sec:privacy-leakage-channels}节）考察系统暴露的信息。这些现象本身不要求存在攻击者。这里额外规定攻击者按其知识、权限和资源选择变化，并追踪变化如何经过系统响应和下游动作破坏安全边界。

\section{对抗安全的基本概念}
\label{sec:security-threats}

\subsection{对抗攻击及其典型表现}
\label{sec:security-attack-definition}

攻击者可以改写系统读取的内容，也可以操纵训练材料、模型组件或交互过程，以破坏预期任务或安全要求。我们把这种有目的的行为称为\term{对抗攻击}（adversarial attack）。本章的攻击目标既包括错误分类，也包括取得受限信息、异常占用服务资源和诱导越权执行；小范围输入扰动只是其中一种实现。

判定攻击首先要看到有目的的选择。偶然误判一封邮件给出了错误样本；反复试探并提交能逃过过滤的写法，则利用了错误来实现目标。还需另判攻击是否成功：系统可能拒绝输入，执行器可能拦下错误提案，因而发起攻击并不必然造成损害。

攻击者能够施加影响的位置称为攻击入口。我们沿邮件系统依次观察输入、训练材料、查询结果和外部内容怎样被利用。下面的情境用于解释机制，均为教学构造。

部署后的模型若已固定，攻击者仍可在推理时操纵输入或交互，使判断偏向自己的目标；这就是\term{规避攻击}（evasion attack）。其中常见的\term{对抗样本}（adversarial example）在预先规定的修改范围内引出不期望的预测。修改范围可以是数值扰动、文本改写或任务允许的物理变化，不能一概限于小范数。我们还要核实修改是否保持任务标签，防止把任务已改变的输入计为攻击成功。比如改写钓鱼邮件措辞后，邮件仍以欺骗为目的，却被分类器判为正常。

攻击也可以先改变学习过程。\term{数据投毒}（data poisoning）操纵训练数据，使训练得到不同模型；\term{后门}（backdoor）则指模型在特定触发条件下出现异常，正常输入上的性能可能仍然良好。因此，干净测试集的准确率不足以排除后门。协作或联邦训练中的恶意更新也是训练入口，需声明恶意参与者比例和可控内容。邮件系统中，攻击者可以污染反馈，让某类钓鱼邮件被学为正常，或让带指定标记的邮件部署后被放行。投毒描述植入操作，后门描述条件性行为；后门既能通过数据植入，也能随受污染参数进入系统。

图~\ref{fig:backdoor-poisoning}将植入过程与部署触发分开画出：同一模型被调用两次，只有触发条件不同。标记并未把钓鱼邮件变成正常邮件。

\begin{figure*}[htbp]
  \centering
  \includegraphics[width=164.600mm]{figures/02-foundations/03-trustworthiness/generated/backdoor-poisoning/figure.png}
  \caption{投毒操作与后门触发分别位于训练和部署。污染记录或参数经训练入口影响过滤器；同一钓鱼邮件无标记时隔离，加触发标记后放行，两个过滤器图案为同一参数模型的两次调用。图中参数污染经训练接入只表示此处路径，并非所有后门都需数据投毒；图未给成功率。}
  \label{fig:backdoor-poisoning}
\end{figure*}

查询接口给出了另一种入口。模型提取（model extraction，也称模型窃取）以恢复参数或结构、或在预定评价分布上复制输入输出功能为目标；恢复私人训练内容另归第~\ref{sec:privacy-concepts}节的信息泄露。两种模型提取目标分别检查参数恢复误差和独立输入上的功能一致率，并报告查询次数、输出粒度、先验知识及替代模型训练成本。例如，在预算内收集分类器回答，再训练替代模型，最后用独立任务输入比较两个服务的判断。

检索和工具使用使系统还会读取第三方内容。\term{提示注入}（prompt injection）把攻击要求混入这些待读内容，试图改变原任务；内容来自网页、文档或工具返回值时，称为间接注入。\term{越狱}（jailbreak）的目标是突破既定行为限制。二者能够共同出现，但应分别声明内容来源、攻击目标和系统权限。比如摘要助手读到邮件中的“归档其他邮件”要求，攻击利用的是内容角色与操作权限的处理，不需要肉眼不可见的扰动。

\subsection{攻击影响与损害路径}
\label{sec:security-damage-paths}

评价之前必须把正常任务和攻击目标分开。逃过检测、诱导指定错误、提取训练信息、复制模型功能及执行未授权动作，各自需要不同成功事件。若统统计入一个攻击成功率，数字便无法说明损害发生在哪个环节，也无法指引防御位置。

同样的分类变化，对单封邮件与整个服务可能有不同影响。放行钓鱼邮件破坏决策完整性；大量隔离正常邮件可能使用户放弃过滤，损害可用性；泄露其他邮件涉及机密性。考虑一个只在“正常邮件”分数超过阈值时放行的分类器：提高阈值可能减少钓鱼邮件被放行，也可能提高正常邮件隔离率。按三类损害分别观察，才能同时报告安全收益与服务代价。

用户只授权归档某类邮件时，邮件正文中的“立即归档全部邮件”仍然是待分析内容，不会增加授权。若系统据此操作其他邮件，关键失效是把内容当作了行动依据。模型复述这段话，与实际越权归档需要分开判定。

我们沿“可控入口—系统响应—下游动作—损害”记录一次攻击。输入导致预测翻转、污染使学习规则改变、工具返回值诱导越权提案，分别是路径中的中间事件。最终损害必须在对应环节核验：出现越权提案证明行为受到操纵，执行端允许并完成该操作才造成行动损害。

入口与损害并非一一对应。改写当前邮件和提前污染训练数据都能以放行钓鱼邮件为目标，但需要不同权限、预算和试验流程。评价因而同时说明攻击从哪里进入、试图破坏哪项要求，再使用相应成功判据。

可用性还可能在答案保持正确时被破坏。受控测试可以观察持续占用计算或重复调用是否拉长队列等待、降低正常请求完成率并增加资源开销。单位请求和整个攻击过程的成本都应报告，并与同负载的正常请求比较；否则容量不足可能被误认成攻击机制。

\subsection{攻击能力与威胁模型}
\label{sec:security-capabilities}

为了让测试可复现，我们用\term{威胁模型}（threat model）固定攻击目标、可观察信息、可修改对象和资源限制。至少应明确受保护资产与安全性质、观察及修改接口、预算和成功事件。可信计算基、初始位置、取得能力的时序与持久性、身份及串通范围，按系统和攻击面继续记录；没有非平凡取值时说明不适用。入口回答从哪里操纵，知识和权限回答如何操纵，目标回答什么结果算成功。这些条件可以同时成立，不能列成互斥攻击类别。

安全要求须放在这个条件范围内：允许的攻击过程不产生指定损害，或损害风险低于容许水平。这里的“允许”是测试赋予的能力，不是现实行动许可。比如允许改写邮件而不允许直接写收件箱数据库，才能检验模型和执行机制的防护；把后者错误纳入同一测试会绕过原本要评价的环节。

知识条件决定攻击者能利用什么信息。白盒提供参数或梯度，黑盒可能仅返回查询结果；结构、版本和错误信息也可能部分可见。权限还需单列：读取参数有助于搜索输入，修改参数却能直接改变被测模型，两者不能互推。

预算则限制查询次数、改动幅度、污染数据量和可控制参与者数。在更强权限下发现攻击，可以提示机制缺口，但未必能在较弱权限下复现。较小预算内搜索失败也只支持这一预算下的测试结论，不能外推到更大预算。

白盒和黑盒只说明知识范围，不评价攻击动机或现实程度。设计者使用完整模型做白盒搜索，是为了检查防御是否仅靠攻击者暂不了解系统；面向公开服务还需按实际反馈做黑盒测试。类别、概率和内部表示提供的信息不同，试验须真正使用已经赋予的权限。声称白盒测试却只运行不充分利用梯度的简单搜索，不能支持相应覆盖结论。

图~\ref{fig:trust-threat-model-coordinates}固定一次黑盒邮件测试的可见信息、可控对象和成功事件。类别反馈不赋予参数修改权，邮件改写权不赋予数据库写权限；查询预算还要另外声明。

\begin{figure*}[!htb]
  \centering
  \includegraphics[width=164.600mm]{figures/02-foundations/03-trustworthiness/tikz/trust-threat-model-coordinates.pdf}
  \caption{一次教学威胁模型的压缩视图：仅返回预测类别，可在欺骗目的保持的前提下改写邮件，以钓鱼邮件判正常为成功事件。20次查询是预算而非实验结果。可见信息、可控对象和成功事件不能互推；完整记录还含资产、信任边界、时序、身份与串通范围，无适用取值的字段明确说明。}
  \label{fig:trust-threat-model-coordinates}
\end{figure*}

图~\ref{fig:trust-security-entrypoints}进一步按系统位置标出入口。位置确定后，仍须组合可见信息、修改权限和成功判据，才能得到具体威胁模型。

\begin{figure*}[htbp]
  \centering
  \includegraphics[width=164.600mm]{figures/02-foundations/03-trustworthiness/tikz/security-entrypoints.pdf}
  \caption{攻击入口按可操纵对象定位：记录内容、更新参数、发布组件、接口输入及查询、提案的行动对象。紫色箭头标控制位置，未表示攻击成功；灰色路径标系统阶段联系。同一攻击可经过多个入口，知识、预算和损害目标仍需组合声明。}
  \label{fig:trust-security-entrypoints}
\end{figure*}

\section{脆弱性的评估}
\label{sec:security-evaluation}

给定攻击能力和预算后，脆弱性评估寻找能够造成损害的路径。为把不同检查放在同一口径上，我们先确定维度、指标和试验程序，再分别检查推理输入、受污染训练、模型与组件供应链、内部机制、查询接口及带权限状态的完整任务。输入试验、内部分析和系统审查可以共同提供证据。

六类对象可串在同一损害路径上：训练材料中的后门进入参数，部署输入触发异常，再经工具执行造成损失。分段测试负责定位薄弱处，系统结论则要核对整条路径。白盒搜索、黑盒试探、组件核查和形式化验证是取证方法，应按被测对象能提供的信息及所需性质选择。

\subsection{评估维度与判定标准}
\label{sec:security-risk-assessment}

一个评估设置至少由三个维度组成。被测对象说明检查输入、内部机制还是工具链；能力和预算说明白盒或黑盒、查询次数及污染比例；结果说明攻击成功、损害及正常任务表现。同一输入入口可以分别在白盒、黑盒和多个预算下测试，因而这些维度需要组合声明。

结果先回答多少任务被攻破，再说明所需查询或修改、成功后的影响，以及正常任务付出的代价。攻击成功率、预算曲线与污染强度、状态变化与恢复成本、干净任务效用与误拒和延迟，分别服务这四个问题。覆盖范围还限定数字对应的攻击、数据和系统配置；未经说明，不能把这些量压成一个安全分数。

\paragraph{可重复的评估程序}

攻击试验按准备、搜索、执行核验和比较展开。准备固定测试者得到的信息，搜索产生候选，执行核验同时检查权限预算与实际系统结果，比较则分开目标达成、损害和正常任务代价。后续试验复用这条程序，但材料、搜索方式和核验对象随入口改变。内部核查和形式化验证也可直接检查计算结构或约束，不必先构造攻击候选，应按实际证据范围报告。

\begin{enumerate}
 \item 冻结模型或训练算法版本、预处理、权限配置及正常任务，保存无攻击时的输出或运行状态；将入口、可改对象、可见反馈、预算和成功事件写成可核验条件。
 \item 测试攻击者只接收约定信息，并在预算内生成候选。保存每次查询、反馈、搜索设置和最终提交对象；预算改变时另列试验设置。
 \item 在完整系统中执行有效候选，核对修改范围和最终输出或状态。随机系统重复整个过程，失败、中止及拒绝运行均予保留。
 \item 按预定试验单位汇总成功与损失，并与正常任务比较。防御改变后先复测已知攻击，再让攻击者在同预算内重新搜索，检查适应防御后的结果。
\end{enumerate}

假定一轮允许改写一封邮件并查询二十次，统计单位便是这一预算内的完整搜索，不能把二十次调用当成二十个独立任务。越权归档的成功事件是授权集合外邮件的状态实际改变；模型产生越权建议另记为行为偏离。

固定$n$次完整试验，令$s_i\in\{0,1\}$指示第$i$次是否成功，经验攻击成功率（attack success rate, ASR）为$\widehat{\mathrm{ASR}}=n^{-1}\sum_i s_i$。分母所含试验必须预先明确。若指标针对原本可正确处理的任务，先冻结这组任务，再计算攻击后失败比例，避免把正常错误算成攻击收益。定向攻击还要核验指定结果，而非任意错误。

\paragraph{攻击结果、损害与证据范围}

成功率还需配套试验数及估计误差。独立同分布的完整试验可使用二项比例区间；同一任务上的相关搜索则保留任务作为统计单位，查询重复不能增加独立样本量。构造攻击、选择攻击和最终比较应使用职责明确的任务集合，避免只报告构造阶段已经攻破的样本。

损失量补充成功事件的严重程度：一次误分类与一次大范围泄露不必有相同代价。令随机元素$\Xi$包含完整试验中的任务实例、环境状态、系统随机性和攻击策略自身的随机数。允许集合$\mathcal A$中的$a$是在抽取$\Xi$前固定的策略规则，只能读取威胁模型开放的信息。执行时，规则可按获准输入和反馈历史决定下一步动作；自适应由这个信息接口体现，不能看完试验结果再逐结果选择策略。记损失为$L(a,\Xi)$，规定能力下的最不利风险写为
\begin{equation}
 R_{\mathrm{adv}}=\sup_{a\in\mathcal A}
 \mathbb E_{\Xi}[L(a,\Xi)]
 \label{eq:security-adversarial-risk}
\end{equation}
上确界比较允许策略的期望损失。只准预先选固定动作时，$\mathcal A$应限于常值策略；若可观察当前输入或先前反馈，则把这些观察写入策略接口。这通常改变风险，不能在$\mathcal A$未变时把上确界移入期望。有限测试只覆盖部分允许策略，提供已发现风险的证据，不能精确估计未知攻击的风险。

式~\eqref{eq:security-adversarial-risk}采用可补偿损失的期望口径。一次泄露或越权就可能造成严重后果时，还要观察事件概率或尾部风险。机会约束限制违规事件的概率，CVaR汇总最差尾部的平均损失；两者改变了对同一损失分布的评价（可参考\BookSectionRef{sec:decision-risk-criteria-robustness}{第一卷“策略的价值与优化”章的“策略的轨迹价值与风险”一节}）。明确禁止的状态变化则需执行端硬约束；在离散时间下，执行规则若使允许集合内每个状态的所有可能下一状态仍在该集合，就能从允许集合内的初态出发，沿轨迹保持安全状态不变量（可参考\BookSectionRef{sec:control-safety-cbf}{第一卷“策略的实现与轨迹控制”章的“策略的约束与安全”一节}）。较低的最坏期望损失不能替代不可补偿的边界。

比较防御时还要固定正常任务和资源条件，报告影响范围、恢复成本及任务效用。拒绝所有输入可能阻断许多攻击，却也取消了服务。有效比较应指出减少的损害，以及随之增加的误拒和延迟。

预算改变了整个过程的成功机会。若每次尝试独立且成功概率固定为$p$，$k$次查询中至少一次成功的概率为$1-(1-p)^k$，随$k$增加。攻击者利用失败反馈改写下一次输入时，独立性和固定成功概率不再成立；此时直接把规定预算内的完整过程作为试验单位。一次调用的失败比例无法代表允许持续试探的服务。

防御本身也会改变策略。过滤关键词后，攻击者可改写措辞；加入监控器后，监控器又可能成为被操纵对象。测试应提供与威胁模型相符的防御知识或试探机会，再比较新搜索结果。旧攻击集用于回归已知问题，不能独自界定改后系统的安全边界。

若需要覆盖规定集合的全部元素，就要补充形式化认证：在输入集合或程序条件下证明性质。证明覆盖的是所写集合，实际系统能否套用仍取决于模型、输入空间和权限假设。比如第~\ref{sec:robustness-certification}节的局部变化界限定数值扰动，不覆盖训练投毒或工具越权。

\subsection{推理输入的评估}
\label{sec:security-prediction-test}

普通测试检查给定输入的预测，对抗测试还在允许修改范围内寻找反例：标签应保持，模型判断却出错。Szegedy等以输入优化展示神经网络的局部脆弱性，并观察到部分扰动在网络间迁移。这说明独立测试集上的分类正确，尚未覆盖样本周围所有允许输入的决策行为\scite{szegedy2014intriguing}。

搜索前需取得原始输入与真实标签、完整预测接口和允许修改集合。白盒额外提供求梯度的模型信息，黑盒则按约定仅返回类别或概率。测试固定模型而改变输入，先核验预算、合法性及标签保持，再在真实接口上确认候选的预测。

\paragraph{利用输入与输出关系搜索反例}

线性打分已能说明小改动如何累积。设$s(\bm{x})=\bm{w}^{\top}\bm{x}+b$，$\bm{x},\bm{w}\in\mathbb R^d$，每个坐标最多改动$\epsilon$。选$\delta_j=\epsilon\operatorname{sign}(w_j)$，各项贡献同向相加，分数增加$\epsilon\|\bm{w}\|_1$，可能跨越分类边界。这个计算只展示一种机制：单特征变化小，不保证总分变化小；它没有断言所有模型都失败。

要把正分数降到负侧，可沿反方向修改。取教学评分器$s(\bm{x})=x_1-x_2+0.05$，正分数判类别1。原输入$(0.5,0.5)$得0.05；每坐标预算0.1时，$(0.4,0.6)$得$-0.15$，预测翻转。该候选只有在任务确认邻域真实标签仍为1时才有效；若真实标签已经变化，不能把翻转计作攻击成功。

图~\ref{fig:trust-adversarial-perturbation}显示这个二维构造的允许区域和候选方向。穿过模型边界说明预测翻转，合法性与任务标签保持则仍需独立核验。

\begin{figure}[!htb]
  \centering
  \includegraphics[width=\linewidth]{figures/02-foundations/03-trustworthiness/tikz/trust-adversarial-perturbation.pdf}
  \caption{教学线性评分器中的预算内候选。$(0.5,0.5)$得0.05，$(0.4,0.6)$用扰动$(-0.1,0.1)$、无穷范数0.1得到$-0.15$。浅色为模型预测区，实线为$s=0$，虚线方形为允许集合。翻转只对应模型；邻域真实标签仍为1且候选合法时，才是有效攻击反例。单步构造不保证任意模型的攻击成功。}
  \label{fig:trust-adversarial-perturbation}
\end{figure}

将上述逐坐标计算用于损失的局部线性近似，便得到Goodfellow等提出的\term{快速梯度符号法}（fast gradient sign method, FGSM）。记$\bm{g}=\nabla_{\bm{x}}\ell(f(\bm{x}),y)$，在无穷范数预算内求
\begin{equation}
 \begin{aligned}
 \max_{\|\bm{\delta}\|_\infty\leq\epsilon}
 \bm{g}^{\top}\bm{\delta}
 &=\epsilon\|\bm{g}\|_1,\\
 \bm{\delta}_{\mathrm{FGSM}}
 &=\epsilon\operatorname{sign}(\bm{g}).
 \end{aligned}
 \label{eq:security-fgsm}
\end{equation}
每一坐标取与梯度同号的端点，就最大化线性化目标；实际候选还要满足像素范围等合法条件。这只需一次反向传播，但线性化的最优扰动不必是原非线性损失的全局最优解\scite{goodfellow2015adversarial}。

非线性损失需要进一步搜索。沿用第~\ref{sec:robustness-change-scope}节的允许集合$U(\bm x)$，把寻找高损失输入写为
\[
 \max_{\bm{x}'\in U(\bm x)}\ell(f(\bm{x}'),y).
\]
$U(\bm x)$若是连续数值集合且可计算投影，就沿增大损失的方向反复更新，再投影回允许集合。搜索得到的有效候选是反例；有限步未找到候选，并未认证集合内不存在更强扰动。

多步搜索在每次更新后重算梯度，补充FGSM的一次近似。这里限定连续数值输入，取
$U_\infty(\bm{x})=\{\bm{x}'\in\mathcal X:\lVert\bm{x}'-\bm{x}\rVert_\infty\leq\epsilon\}$，$\mathcal X$为合法输入盒，记$\bm{g}^{(t)}=\nabla_{\bm{x}^{(t)}}\ell(f(\bm{x}^{(t)}),y)$。常用更新为
\begin{equation}
 \bm{x}^{(t+1)}
 =\Pi_{U_\infty(\bm{x})}
 \left[
 \bm{x}^{(t)}
 +\eta\operatorname{sign}(\bm{g}^{(t)})
 \right],
 \label{eq:security-pgd}
\end{equation}
$\Pi_{U_\infty(\bm{x})}$是到闭凸允许集合的欧氏投影，轴对齐盒中可逐坐标裁剪；$\eta$为步长，初值可在允许区内随机抽取。这类攻击称为PGD，借用投影梯度下降（projected gradient descent）的缩写，但此处沿增大损失的方向运行。文本等离散集合须另定搜索与可行化方法。随机重启探索不同区域，更多迭代细化一条路径，二者不能互代；某个设置的损失停止改善，也不足以确认最坏输入\scite{madry2018adversarial}。

\paragraph{针对防御机制调整搜索}

搜索还要适配防御。除不同初始化和策略外，必要时补充迁移攻击；随机或不可微处理则需要相应搜索机制。若处理仅令梯度难用而决策边界仍脆弱，就发生\term{梯度掩蔽}（gradient masking）。原梯度攻击失效不能独自证明防御有效。

目标函数也会使搜索受阻。Carlini与Wagner评估防御蒸馏时改用分类前的打分差，避免输出概率饱和。固定像素已缩放到$[0,1]$的图像输入$\bm{x}\in[0,1]^d$，像素盒只保证数值合法，标签和语义仍另验；针对类别$k$可采用
\begin{equation}
 \begin{aligned}
 \min_{\bm{x}'\in[0,1]^d}\quad&
 \|\bm{x}'-\bm{x}\|_2^2+c\,h_k(\bm{x}'),\\
 h_k(\bm{x}')={}&
 \max\left\{\max_{j\ne k}
 [z_j(\bm{x}')-z_k(\bm{x}')],-\kappa\right\},
 \end{aligned}
 \label{eq:security-cw}
\end{equation}
$z_j$为类别打分，$c>0$控制软惩罚，$\kappa\geq0$指定希望达到的间隔。有限$c$的优化结果未必成功：$h_k(\bm{x}')\leq0$仅说明目标打分不低于其他类别，$h_k(\bm{x}')=-\kappa$才说明至少领先$\kappa$。实际需搜索$c$，再用真实分类器及平分规则确认目标输出。数值搜索困难不是防御证据，合适目标可能重新暴露反例\scite{carlini2017evaluating}。

不可微处理同样可能隐藏可用方向。Athalye等分析量化和离散操作如何阻碍梯度却保留脆弱边界：反向传播近似在前向使用真实防御，在反向用可微近似搜索；随机防御可多次抽取变换，平均损失来估计期望梯度。候选最终必须在真实防御上成功，近似模型中的成功不算实际结果\scite{athalye2018obfuscated}。

这些选择改变的环节不同。FGSM与PGD改变更新次数，CW改变搜索目标及扰动代价，反向近似与随机变换平均改变取方向的方法。因此，同一试验能够组合多步更新、打分差目标和变换平均。组合依据是模型接口与防御机制，判定依据始终是候选在真实系统中的结果。

\paragraph{互补搜索与实际输入条件}

单一搜索可能漏掉某种模型或损失的弱点。AutoAttack组合自动调步长的梯度搜索、边界搜索和无梯度查询搜索，以互补机制减少手工调参及单一方法造成的漏检。“免参数调节”指流程预设相关规则，预算仍需声明，方法对具体防御的适用性仍需检查\scite{croce2020autoattack}。

用教学数据核对组合口径：100个样本上两个攻击分别发现12和15个失败，重合8个。联合发现$12+15-8=19$个，不能平均成13.5个；这组搜索测得的鲁棒准确率为$1-19/100=81\%$。其余样本只是尚未找到反例。在相同样本与允许集合下，真正最坏情况的鲁棒准确率至多为81\%。

物理入口又改变了输入条件。提交图案经过拍摄、缩放和压缩后，未必保持数字数组中的攻击效果；只直接输入数组的试验也只覆盖这一入口。若评价物理攻击，需记录成像距离、角度和光照；以随机变换描述它们时，还要说明变换分布如何代表实际条件。

汇总时先固定集合和预算，再在原本预测正确的样本上按扰动强度计算成功率。指定类别输出才是定向成功，任意错误才是非定向成功。同一个样本接受多种搜索时，按是否至少出现一个有效反例汇总，查询不能当作独立样本。

\subsection{训练过程的评估}
\label{sec:security-training-test}

训练入口改变了被测过程。输入攻击运行固定模型即可；投毒试验则要把受污染材料交给学习者，再检查训练结果。攻击者如何选污染内容，与学习者采用哪些配置，均须固定。

对照使用相同训练流程：原数据一份，另一份只按规定入口加入或修改污染。采用配对初始化和超参数，在多个随机种子下重复，留出集既不参与污染构造也不选最有利防御。若知识仅覆盖替代模型，污染先在替代流程中生成，再检验向真实训练迁移；不能暗用真实梯度改变知识条件。

为描述训练响应，把正常数据记成序列$D=((\bm{x}_i,y_i))_{i=1}^{n}$，保留索引和重复项。$\mathcal P(D)$收集规定追加、按索引替换等操作及预算下的完整受污染序列，$D'\in\mathcal P(D)$不使用会折叠重复记录的集合并集。先限定白盒设定：攻击者知道$D$、真实算法、目标损失和种子分布，却在提交$D'$时不知道本轮实际种子。学习者用固定算法$\mathcal T(D';\xi)$响应，$\xi\sim Q$含按预定分布抽取的初始化、数据顺序及其他随机性。相应定向投毒目标为
\begin{equation}
 \begin{aligned}
 \max_{D'\in\mathcal P(D)}\quad&
 \mathbb E_{\xi\sim Q}
 \left[L_{\mathrm{target}}\bigl(\bm{\theta}_{\mathcal T}
       (D',\xi)\bigr)\right],\\
 \text{其中}\quad&
 \bm{\theta}_{\mathcal T}(D',\xi)
 =\mathcal T(D';\xi).
 \end{aligned}
 \label{eq:security-poison-bilevel}
\end{equation}
$L_{\mathrm{target}}$按攻击目标选择，值越大表示损害目标越接近，不必等于正常训练损失。外层选数据，$\mathcal T$给出学习者响应，并须明确优化器、超参数、停止和多解选择；确定性训练可省去$\xi$期望。种子若对攻击者可见，或评价改用最坏种子及分位数，需相应调整信息条件和统计量。实验报告多个种子的分布，不能用多值$\argmin$冒充确定响应。只知替代流程时仍先在替代流程构造$D'$，真实$\mathcal T$只负责迁移评价，不能直接套用白盒目标。这里污染经学习影响后续多次预测，与只改变一次输入不同。

BadNets展示外包训练和模型供应链的后门风险：让触发样本指向目标类别，同时保持正常输入表现。普通验证集若不含触发条件，就可能看不到异常。后门还可随预训练模型迁入下游，因而下游数据来源检查不够覆盖模型入口\scite{gu2017badnets}。

人工看来正确的标签也挡不住所有投毒。Poison Frogs的\term{干净标签投毒}（clean-label poisoning）保持合理标签，却让污染样本在表示中靠近指定目标，移动分类边界。这说明标签复核不能替代训练安全审查。效果依赖表示、训练方式及目标知识，不能外推为任意少量数据控制任意模型\scite{shafahi2018poison}。

两类构造虽都改训练材料，目标范围不同。BadNets式后门学习触发与目标类别的联系，让多种输入加触发后同向偏离；干净标签投毒可针对单个目标，以表面标签正确的样本改变表示附近的训练约束。前者通常验收一组触发输入，后者还核验预先指定目标。构造与判据均应匹配数据权限及目标范围。

污染收益必须与无污染训练比较。正常模型已常误判目标时，不能把全部污染后错误归因于攻击。匹配配置下报告目标错误率、干净性能和初始化变化，并明确污染比例及记录是否实际进入训练。后门同时检查触发时目标成功和无触发性能，分别显示有效性与隐蔽性。

固定类别$k$和真实标签非$k$的留出输入，每个保留无触发与触发版本。配对转移率的主分母取污染模型在无触发版不输出$k$的样本，再计触发后转为$k$的比例；更严格时再要求原版分类正确。另行报告全部非目标样本的触发ASR、无触发错误率及未污染模型在同一配对输入上的结果，以分开触发转移与基线错误。跨防御时冻结评价集合，按污染预算和触发强度展开指标，不能只挑最有效的一次训练。

来源、标签和表示异常给出材料线索，受控投毒才直接检查材料进入学习后的影响，二者应配合。比较网络结构时，尽量匹配数据、算力和权限，并让攻击各自适应结构；正常准确率不能单独回答投毒敏感性。已交付模型若不能重跑原训练，则转向内部诊断，并记录缺少训练资料的证据限制。

\subsection{模型与运行组件的供应链评估}
\label{sec:security-supply-chain-test}

交付模型还会引入序列化格式、预处理代码、运行依赖和加载权限。供应链检查先冻结版本清单，核验签名、摘要及来源，再隔离记录加载、初始化、推理和外部访问。文件替换、依赖变更、预处理旁路，以及合法来源的恶意行为均需覆盖；签名不排除后一种情形。来源真实、制品完整和行为安全是分别验收的结论。

隔离运行尽量缩小文件、网络、凭据和执行权限，同时检查加载是否执行未声明代码、访问清单外资源，预处理是否符合冻结版本。异常定位可疑组件，未见异常也只覆盖当前路径与环境。供应链证据继续与内部及完整任务测试合并，才能核实交付物是否构成损害路径。

\subsection{模型内部机制的评估}
\label{sec:security-model-test}

内部审查使用计算图、参数、梯度或中间表示，寻找机制或验证行为性质，前提是具备相应访问权限。只有查询权时应报告黑盒结果的范围。审查者获得的内部权限还须与真实攻击者的权限分开声明。

架构名称、层数及参数量不能单独决定安全；同一结构的不同数据和参数可产生不同边界或后门。输入测试提供失败候选，内部分析则解释失败路径、检验搜索充分性并尝试在指定条件下排除反例。单列内部评估，是为了明确这部分新增机制与验证证据。

\paragraph{计算路径与局部敏感性的诊断}

先确认预处理、网络各层、随机分支和输出变换都进入实际被测路径，再查看梯度和激活。单点梯度能指引搜索，小梯度却可能只是饱和或掩蔽，不能作邻域变化上界。表示异常也只是候选线索，须回到预测判断是否实现攻击目标。

机制怀疑可以用受控组件比较检验：在副本中替换或移除组件，其他条件保持一致，再测攻击和正常任务。比较量化前后的真实失败与梯度可用性，就能帮助区分消除脆弱性和阻碍搜索。副本结果仍不能冒充原模型结论；第~\ref{chap:trust-interpretability}章的过程解释给候选机制，安全判断再补有效攻击或性质验证。

\paragraph{异常触发机制的诊断}

已有模型的后门诊断可从“是否有一种触发将多类输入推向同一目标”入手。Neural Cleanse逐目标类别反推触发，用掩码$\ell_1$范数等量度刻画大小，再以类别间中位数及中位绝对偏差寻找异常小的触发。异常提示需继续查样本、过滤或修复；它把不可穷举的搜索变成可操作诊断，却依赖触发结构和异常性假设。输入相关触发、多受影响类别及自然类别不对称都可能削弱判据\scite{wang2019neuralcleanse}。

诊断候选还不证明真实后门存在。应在独立输入上核验目标类别比例及触发前后变化，检查触发的可实现条件；若怀疑投毒，还追溯训练和模型来源。没有候选，只说明本次搜索及结构未找到异常。线索如何用于来源检查和参数修复，见第~\ref{sec:security-training-defense}节与第~\ref{sec:security-model-defense}节。

\paragraph{利用计算结构验证行为约束}

内部计算还可转为约束求解：同时规定允许输入、网络计算和不希望出现的输出，寻找满足全部条件的反例。ReLU网络的线性层给线性等式，激活与不激活给分段约束。Reluplex将这些计算规则纳入性质验证，覆盖的对象超出了有限候选输入试验\scite{katz2017reluplex}。

设原输入$\bm{x}$的正确类别为$y$，类别$j$的分数为$z_j(\bm{x}')$，允许集合为$U(\bm{x})$。检查是否有$\bm{x}'\in U(\bm{x})$及$j\ne y$满足$z_j(\bm{x}')\geq z_y(\bm{x}')$，就是搜索至少追平正确类的竞争类。核验解后仍需按实际平分规则确认误分类；可靠求解器证明不可满足时，才排除集合内全部这类候选。超时和未知结果保留为未判定。

也可沿网络传播有效上下界。每个竞争类别在整个集合内的$z_y-z_j$下界严格为正，才能认证分类不变；非正下界可能过松，并未给出攻击。条件及报告方式见第~\ref{sec:robustness-certification}节。结构验证针对给定参数和性质，范围不继承到训练投毒、语义改写或工具权限。

\subsection{查询接口的评估}
\label{sec:security-information-test}

查询也能以复制服务功能为目标，而不制造单次错误。Tramèr等指出，预测概率等辅助信息能提供参数方程。以已知二分类逻辑回归为例，输入$\bm{x}\in\mathbb R^d$、权重$\bm{w}\in\mathbb R^d$和截距$b$，返回类别1的精确概率$p(\bm{x})$，则
\begin{equation}
 \log\frac{p(\bm{x})}{1-p(\bm{x})}
 =\bm{w}^{\top}\bm{x}+b.
 \label{eq:security-model-extraction}
\end{equation}
预处理已知、概率在零与一之间且查询设计满秩时，$d+1$个独立方程确定$d+1$个未知量。舍入、隐藏特征和复杂模型改变实际成本；仅类别反馈虽不再给等式，仍可用于训练替代模型\scite{tramer2016stealing}。

一维教学例中，$x=0$返回$p=0.5$，$x=1$返回$p=e/(1+e)$。由式~\eqref{eq:security-model-extraction}得$b=0$、$w+b=1$，故$w=1$。再用未参与求解的输入核验输出，检查接口假设。可解性的关键是精确概率转为对数几率后的线性方程，而非查询的高置信本身。

类别反馈无法计算上述对数几率，可以改做替代训练：提交允许输入，收集标签，训练替代分类器，再用剩余预算补充当前模型容易分歧的区域。方程求解与替代训练都能复制功能，但依赖不同反馈和模型假设。私人内容提取仍按上一章的秘密与泄露判据另验。

参数恢复和功能相似需分别测量。网络权重不同也可预测接近，仅查权重会漏掉功能复制。正常用户常见输入、边界附近和罕见输入的一致率也可能不同，评价分布须声明。预算同时计查询量、每次反馈信息量及替代训练成本。

具体试验让攻击者在规定预算内训练替代模型，用与查询集独立的任务输入测量一致率和任务效用，再画预算变化曲线。个人信息或训练内容目标则换用第~\ref{sec:privacy-verification}节的成员推断或恢复指标，功能一致率不构成泄露证据。收益连同输出粒度、限额和已知信息报告。

接口可用性另用资源试验检验。隔离环境中固定并发、输入规模、超时和总预算，比较正常与可疑请求的开销、完成率和等待时间。正常流量同样触及容量上限时据实报告；少量允许请求引起持续异常消耗时再定位解析、推理或调度路径。复制功能、泄露私人信息和耗尽资源共享入口，成功及损失判据却不同。

\subsection{交互与执行链的评估}
\label{sec:security-interaction-test}

生成系统需把攻击者想要的文本与实际行动分开。Greshake等的间接提示注入研究追踪第三方内容经检索、工具调用进入上下文，如何从数据来源变成行为操纵入口。测试必须覆盖来源和授权边界，仅查用户直接输入的恶意要求不足以覆盖该路径\scite{greshake2023indirect}。

图~\ref{fig:trust-prompt-injection-boundary}分别标出内容来源与行动授权。越权提案能否变成损害，取决于执行链对提案的处理。

\begin{figure*}[!t]
  \centering
  \includegraphics[width=164.600mm]{figures/02-foundations/03-trustworthiness/tikz/prompt-injection-boundary.pdf}
  \caption{第三方邮件影响动作提案的教学情景。用户只授权邮件A，外来内容用虚线夹带删除B要求，模型可能提出归档A及删除B。执行端以可信授权和状态独立检查，允许A而拒绝B，故B不变。这说明防护职责，实际仍需检查旁路、状态变化与损害。}
  \label{fig:trust-prompt-injection-boundary}
\end{figure*}

自动搜索提高了越狱测试强度。Zou等用词元梯度筛选降低攻击代理损失的候选替换，再以前向计算选择离散修改；跨请求和模型联合优化用于研究迁移。代理损失下降仍须最终行为确认，迁移结果也不覆盖所有版本与系统防护\scite{zou2023universal}。

AgentDojo在带工具和状态的环境同时评价用户任务完成与攻击者指定偏离。正常任务指标能区分保留能力的防御和一概拒绝工具结果的做法，固定任务也便于比较注入位置。证据仍受环境、权限和攻击集合限制，不能代替部署授权审计\scite{debenedetti2024agentdojo}。

一轮测试运行完整任务：按规定位置注入，保存读取、动作提案、执行校验和状态变化，分别判断用户任务完成与攻击目标达成。每任务设无注入对照，再改变来源、权限或预算比较。脱离执行的文本试验能测诱导，却不能独立测越权损害。

归档任务可用集合把四种结果分开。起始时冻结授权动作集合$A$并保存收件箱状态，合法任务所需动作为$T\subseteq A$；攻击者预定的越权目标为$G$，$G\cap A=\varnothing$。攻击者仅可修改一封待读邮件正文，系统随后记录提案集合$B$、执行器接受集合$C$及最终状态。$B\not\subseteq A$表示提案越权，$C\not\subseteq A$表示执行端允许越权；合法完成要求最终状态等于只执行$T$的预期状态，攻击成功要求状态实现$G$。四项分计，才能区别未受诱导、诱导被拦截和全部拒绝造成的低成功率。

判定工具也要核验。关键词可能把引用与执行混淆，模型裁判也可能被待评文本操纵。工具任务优先查外部结果，如未授权发送或越权修改。无法自动判断的案例交独立复核，记录标准与不确定性，不用含混分数替代。

\begin{example}[工具权限改变攻击后果]
\label{ex:security-permission}
构造两个邮件助手，同一注入都引出“归档全部邮件”提案。甲可直接调用无范围限制的归档接口；乙只获得本任务授权标识，执行端再复查。同样的文本在甲中能转为越权行动，在乙中受执行权限限制，因而评价分开模型诱导与实际损害。
\end{example}

先将模型与执行器分开，可定位授权规则。隔离环境直接提交正常、含越权对象及超数量请求，核对底层状态，再接入受注入模型重测。前者检验规则，后者找组合旁路；来源标记、路由、持久记忆和重试也入路径清单，范围不能停在单轮回答。

\section{攻击防御与损害控制}
\label{sec:security-defense}

防御承接已定位的损害路径，明确在哪一环节干预。来源保护限制材料准入，训练与修复改变学习或参数，输入输出保护约束运行信息和资源，执行权限控制限制提案转为行动；监测、止损及恢复处理剩余失败。同一攻击链可由多类措施共同阻断，每类都需对应试验。

恶意邮件可能被输入检查拒绝，也可能被模型忽略；两层都失败时，独立执行器仍可拦下越权归档。若内容已经写入持久状态，还要隔离恢复。这一分类按作用位置与对象组织，单独采用对抗训练不能证明来源、查询及权限均安全。

\subsection{数据与模型来源的保护}
\label{sec:security-training-defense}

材料准入先核查数据来源、提交主体、标注变更，以及模型和处理代码版本。数字签名只有在密钥与身份关联可信并验证通过时才支持来源和完整性；文件摘要须比对可信渠道的预期值才能发现变化。单个摘要不认证发布者，来源与完整性通过也不证明材料无投毒或后门。

这一步回答从哪里取得交付物、取得后是否改变。运行行为是否安全，仍需第~\ref{sec:security-supply-chain-test}节的隔离验收。

准入链可独立于预测：登记来源，核验主体及数量，新来源暂存审查，获准对象列入可追踪训练清单。分别以未授权来源、超配额和合法来源污染测试。前两类检查权限规则是否执行，后一类检查内容审查能否发现污染，同时报告合法材料误拒。

合法来源仍可能提供异常内容。标签抽查找错误标注，重复与来源分布找集中提交，表示分析找表面标签正常的污染。各检查只针对相应线索，需搭配受控污染试验和合法样本误拒评估。

谱特征方法（Spectral Signatures）借候选模型的类别内表示筛查：中心化表示，求变化最大的方向，以投影大小选候选污染，清理后重训。Tran等利用所研究后门在表示中的可分离结构，条件包含污染比例与表示分离；不能据此认定所有后门都有相同异常\scite{tran2018spectral}。

表示筛查需要候选训练，宜在隔离环境完成训练与清理，再决定材料能否用于后续学习。它形成借模型改善数据质量的循环，入库字段不足以独立完成；审查用途也不赋予候选模型生产准入。

第~\ref{sec:security-model-test}节的诊断帮助挑选模型、样本及触发线索，但不会自行清除异常。修复要记录改了哪些材料、参数或输入规则，再重测触发成功率。

异常只是调查线索，不能直接作为批量删除理由。稀有合法群体也可能有异常表示，需考虑第~\ref{chap:trust-fairness}章的代表性问题。修复同时检查触发是否消失、正常类别和稀有群体性能，避免把代价转嫁给特定用户。

\subsection{模型训练与修复}
\label{sec:security-prediction-defense}
\label{sec:security-model-defense}

模型防御直接改变学习规则或参数。按抑制对象，下面分别讨论对抗训练降低输入攻击收益、鲁棒聚合限制恶意更新、已交付模型修复削弱异常机制。三者可组合，但输入训练不自动解决投毒，后门修复也需来源审查。

\paragraph{针对输入操纵的训练}

输入防御可把当前搜索到的高损失样本用于训练，近似最小化允许集合内的最坏损失。Madry等用鲁棒优化连接这一目标和攻击搜索。所得保护仍限定于允许扰动、搜索强度与训练分布\scite{madry2018adversarial}。

每个小批次先固定参数，按第~\ref{sec:security-prediction-test}节搜索允许区域内的高损失输入，再固定候选更新参数。记参数$\bm\theta$、模型$f_{\bm\theta}$，训练样本$(\bm x,y)\sim\mathcal P_{\mathrm{tr}}$，要求扰动$\bm\delta$满足$\bm x+\bm\delta\in U(\bm x)$，理想目标为
\begin{equation}
 \min_{\bm\theta}\ \mathbb E_{(\bm x,y)\sim\mathcal P_{\mathrm{tr}}}
 \left[\sup_{\bm\delta:\,\bm x+\bm\delta\in U(\bm x)}
  \ell(f_{\bm\theta}(\bm x+\bm\delta),y)\right].
 \label{eq:security-adversarial-training}
\end{equation}
内层寻找当前模型的薄弱输入，外层减少这些输入的损失。实现通常只有有限步内层搜索，候选也只用于当前更新；参数变化后要重搜，永久复用固定攻击集不等价于该目标。

固定加噪不随模型寻找薄弱点，对抗训练的内层则响应当前参数。但内层不足时，模型可能仅适应训练搜索的模式。最终评估需包含未参与调参的攻击，并比较干净性能与不同预算，不能完全复用训练攻击作为唯一证据。

第~\ref{sec:robustness-local-enhancement}节的TRADES（式~\eqref{eq:robustness-trades}）明确正常预测与局部一致性的权衡，随机平滑半径（式~\eqref{eq:robustness-smoothing-radius}）在条件下提供认证；这里分别将它们用作经验增强和限定集合保证。像素范数认证不覆盖邮件语义改写，输入扰动训练也不覆盖投毒。换攻击面时须重新明确适用假设。

\paragraph{针对恶意更新的训练规则}

协作参与者身份合法，更新仍可能恶意。鲁棒聚合用较少受极端值影响的规则代替平均。Yin等分析逐坐标中位数与截尾均值的分布式学习误差，结论需要相应统计、优化条件，尤其匹配恶意比例与正常更新分布\scite{yin2018byzantine}。

聚合改变参与者向量的合并方式。逐坐标中位数分别选每维中位值。教学一维更新$0.9,1.0,1.1,8,9$中，后两个由攻击者控制，平均为4，中位数为1.1，显示少量极端值对这一坐标的拉动受到限制。高维时各维中位值可能来自不同参与者，正常数据差异也会使更新分散，因而例子不保证聚合向量有用。试验需固定人数及恶意比例，比较攻击后损失、正常收敛和群体损失。

截尾均值先逐维删去规定数量的最大最小值，再平均剩余项；截尾数须匹配威胁条件。大量伪造身份或高度不同的正常数据，会破坏比例与集中性前提。第~\ref{sec:privacy-collaborative-computation}节的安全聚合隐藏个体更新，而鲁棒聚合要执行特定运算；两者组合前必须检查协议能否支持，名称并列不给联合保证。

\paragraph{已训练模型的修复}

有可信后门线索后，可以清理受污染材料并从可信起点重训，或尝试修订交付模型。Fine-Pruning先按干净输入激活剪去部分低活动神经元，再以干净数据微调其余参数。Liu等比较单独剪枝、微调及组合，并研究适应剪枝的攻击，说明修复也要面对攻击者重新选择策略\scite{liu2018finepruning}。

修复需可信诊断和微调数据、正常能力覆盖，以及异常与有用机制可分离。共用表示时剪枝可能损害正常任务，未测条件中也可能残留后门。验收重跑触发搜索、已知攻击和保留任务；消除单个触发的效果不能外推到全部未知触发。

安全修复以限定攻击下异常行为为验收对象，机器遗忘则比较移除指定记录后的训练结果，见第~\ref{sec:privacy-unlearning}节。两者都能改参数，但比较基准不同。依赖查不清或修复不达标时，应改用可信数据及起点重训，不能继续信任受污染检查点。

\subsection{推理输入与查询输出的保护}
\label{sec:security-information-defense}
\label{sec:security-input-defense}

运行接口有两个防护位置。模型前的处理决定哪些内容以何种表示参与计算，交付后的输出和查询控制决定反馈量及尝试次数。这些措施通常不改训练参数，因此应测试包含全部处理的实际接口。

\paragraph{输入检查与可疑请求处置}

格式、类型、长度与范围校验排除非法输入；语义合法但可疑的请求还可由异常、已知触发或独立风险检测决定拒绝、复核或受限处理。拒绝只是处置，不证明攻击识别正确。评价同时计漏检、正常误拒及接受请求的剩余风险。

去噪、压缩、量化或规范化会改变到达模型的输入，也可能损失任务信息或只阻碍梯度。自适应测试应包含处理过程，按第~\ref{sec:security-prediction-test}节用反向近似、变换平均或黑盒搜索检查掩蔽。某种变换的收益不能证明全部语义保持攻击消失。

检索和工具返回值应保留来源，不能直接提升权限或覆盖配置。恶意短语过滤处理已知模式，角色标记帮助区分内容，但改写及未知模式仍要测。外部材料可能同时含合法步骤和恶意要求，最终行动还须下节的独立授权。

\paragraph{查询输出与资源配额的控制}

减少反馈信息、限制查询可以增加模型或私人信息提取成本，同时也损失正常用途。是否需要正式输出保护，应按第~\ref{chap:trust-privacy}章判断；参数不公开并不证明信息无泄露。

反馈粒度直接改变查询约束。式~\eqref{eq:security-model-extraction}的精确概率给参数等式，概率区间给对数几率区间，类别只给边界一侧的信息；查询限额再限制累积量。正常用户也可能失去校准概率或批处理能力，返回内容需按业务用途选择。

固定模型和查询预算，分别用精确概率、有限精度及仅类别训练替代模型，在独立输入上比较功能一致率，再检查正常任务损失。这样才能区分减少信息与普遍拖慢使用。类别查询仍可能逐步逼近边界，报告应展示预算与提取效用关系，而不只标记是否提供概率。

模型保密与隐私保护另有对象。差分隐私限制相邻关系中一个记录或用户的变化如何影响发布分布，却允许发布总体规律，因此功能仍可被近似复制。结构隐藏也不给训练数据保证。接口先决定保护个人信息、商业模型还是安全规则知识，再设对应判据。

改接口后按第~\ref{sec:security-information-test}节复测同预算提取收益，并记录正常限流、精度损失及延迟。多账户和多会话若在攻击权限内，配额须按相同身份边界计量；单会话限额不等于完整攻击预算。

资源配额可限制输入、并发、生成长度、工具次数及任务时长，并隔离租户队列。配额需覆盖完整任务和可识别主体，否则重试或多会话可累积超限。测试同时检查异常消耗与长文本、批处理等合法需求的损失。

\subsection{行动执行中的权限控制}
\label{sec:security-execution-defense}

第~\ref{sec:alignment-intervention-execution}节已把输出限定为动作提案，由独立执行器复核对象、参数和状态。这里进一步检验攻击者能否绕过规则，范围包括可控身份、内容、调用顺序和反馈；仅存在检查步骤不足以支持安全结论。

权限先缩到任务所需：摘要不需删除，读取指定文档不需访问全部私有目录。测试再覆盖伪造标识、批量参数夹带越权对象、检查与执行间状态改变、多账户或连续低权限动作绕过上限。沙箱和结构化参数的作用取决于这些路径受到什么约束，名称不能替代隔离边界。

多层措施还要检查共同弱点。过滤、行为评判和授权若共用同一个可被文本操纵的模型，层数增加未必明显降风险。设计需追踪共享信息，并检查各层失效是否独立。

先直接测试执行器：伪造主体、越权或重复对象、超量批量及状态已变请求；再将同条件置回提示注入的完整交互。前者检验授权规则，后者检查模型、路由及执行器旁路。每轮核对底层状态和正常完成；文本拒绝仅是行为结果，越权对象没有状态变化才支持相应条件下的执行边界。

来源保留同样适用于检索增强。网页建议可以提供任务事实，却不能更改工具授权；引号或“不可信”标记给模型区分信号，仍需模型正确理解。外部检查应在模型忽略信号时仍限制越权。不可逆动作的复核界面需列具体对象和后果，使批准的行动可辨认。

\subsection{攻击监测与损害恢复}
\label{sec:security-recovery}

剩余风险要求监测与响应衔接。异常查询密度、资源消耗、工具范围、连续拒绝、版本变化和外部流向可提供线索，再结合状态及独立记录判断是否需要处置。模型自报未受攻击不能代替证据，误报漏报须演练计量。

响应随传播位置选择：暂停任务和后续工具、撤销会话凭据、隔离数据或版本，必要时切换受限功能或已验版本。外发正在发生时及时关出口，模型污染时停止该版本继续决策；仅拒绝当前输入，不能清除参数或记忆中的污染。

检测可能晚于部分损害。限制动作规模、保存恢复状态、撤权及停止后续操作可缩小影响，能力仍需实测。文件恢复不能收回外发信息，停止后续行动也不能撤销既有不可逆操作。

隔离演练可先保存状态，运行一条错误改邮件标签的教学轨迹，在规定时刻触发检测，检查后续动作是否停止、权限是否撤销及标签是否恢复。记录检测到停止时长、受影响对象、恢复状态及残余损失。已经外发的信息只能检查后续传播是否受限，在本地删除日志不等于收回信息。演练回答发现失效后还能控制什么，不重新替代原攻击判定。

还需保全足以区分模型错误、注入及授权配置问题的证据；日志可能含敏感输入，按第~\ref{chap:trust-privacy}章限制访问与保留。仅替换模型却保留注入入口，或仅删除文档却不检查记忆及派生状态，都可能留下继续触发的路径。恢复范围须沿信息传播确定。

处置后回到第~\ref{sec:security-evaluation}节，同时复测正常任务、已知攻击和适应新防御的搜索。报告明确规定能力预算下受到限制的损害及未覆盖条件。

重新开放还需退出门槛。按冻结清单与传播路径逐项对账受影响资产、依赖及派生状态，无法确认的引用另列未闭合项。恢复版本须来自可信基线或可复核重建；相关凭据、配置、模型、记忆与派生状态完成轮换、清理或隔离。原入口及已知旁路须在规定能力预算下通过回归、自适应复测和底层状态核验，正常任务达到预设最低性能，不可逆残余损失单独登记。硬门槛未过就保持隔离或受限服务，并留可执行回退版本。全部通过也仅支持同一威胁模型下逐步恢复权限，未排除未知攻击。

图~\ref{fig:trust-defense-in-depth}把准入、模型复核、接口、独立授权和监测恢复放在传播路径上。来源检查可以拦截未获准材料，但不能保证发现合法来源污染；格式合法注入也可能穿过输入和模型，授权仍须独立限制对象。恢复负责既有影响，与事前阻断分别记录。

\begin{figure*}[htbp]
  \centering
  \includegraphics[width=164.600mm]{figures/02-foundations/03-trustworthiness/generated/defense-in-depth/figure.png}
  \caption{纵深防御在训练与运行路径上的分工。未授权来源止于准入，异常候选模型暂不发布，但来源通过不证明内容无污染。合法格式的外部注入可到模型并提出归档B，独立授权拒绝越权对象后B不变。监测用独立记录和底层状态，失效后暂停、撤权及恢复可逆状态；外发信息不能据此收回。阻断依赖相应检查触发，未保证全部攻击可发现，也未假定各层失效独立。}
  \label{fig:trust-defense-in-depth}
\end{figure*}

\FloatBarrier
\Needspace{8\baselineskip}
\section{本章小结}
\label{sec:security-summary}

邮件系统的安全判断已经从预测扩展到信息和行动路径。攻击类型描述操纵操作，完整性、机密性和可用性描述损害，威胁模型固定可见、可改及尝试预算。沿这些条件核验时，预测出错、越权提案和实际损害须分别记录，才能看出保护在哪一环节失效或生效。

输入反例、配对训练、供应链隔离、内部机制、查询收益和完整状态各提供一种证据，合起来限定已发现的脆弱范围。单次搜索失败、架构名称及干净成绩不能取代这些检查。成功率只有与预算、损害、正常代价和覆盖范围共同报告，才支持防御比较。

防御选择由薄弱位置决定：来源审查无法自行消除参数后门，训练不能代替配额与授权，过滤不能自行清理持久状态。组合数据、模型、接口和执行保护后，还要开放同条件的新搜索。监测与恢复决定突破后的损失边界，其中外泄信息与可逆状态修改必须分开验收。

第~\ref{chap:trust-alignment}章还讨论无外部攻击时的目标与行为偏离。这里增加有目的的操纵，攻击者可利用预测误差、环境变化、泄露和权限缺口使保护失效。因此，应在同一系统及明确攻击能力内复核前面各章的性质；改善其中一项，仍须检查其余损害路径。

\paragraph{延伸阅读}

Madry等的鲁棒优化与Athalye等的自适应攻击，可用来对照防御训练与评估：前者把扰动范围纳入训练，后者区分梯度搜索受阻与脆弱性消除。继续阅读时，始终分开经验失败和形式化认证\scite{madry2018adversarial,athalye2018obfuscated}。

检索及工具系统可沿间接提示注入与AgentDojo继续阅读：前者追踪第三方内容进入行为链，后者同时判用户任务和攻击目标。比较时保留输出提案、执行端接受动作及最终状态这三个层次\scite{greshake2023indirect,debenedetti2024agentdojo}。
